%! Planning a Study — Unit 3, Lesson 3.1 — Homework (Answer Key)
% MIRROR of ../homework/main.tex: -key for -boxes, the header, and the correct MC option wrapped in
% \ans{}. Every \answerspace and \gradedblank line is identical; key mode prints each display.
% NO teachernote here — teacher prose lives in the lesson plan.
% GRADUAL RELEASE: this is the lesson's SCORED practice and the LAST component in the packet.
% Paper homework is the DEFAULT for this course; DeltaMath is an occasional override that
% replaces this page. Started in the last eight minutes of the period and finished after two
% study halls. The AP Practice component before it is extra credit, scored 1 to 5.
% THIRD CONTEXT: the notes teach on one-line methods, a principal's homework survey and four
% one-line studies; Guided Practice runs on a university study app; AP Practice on reading
% scores, an Ohio park poll and energy drinks. These items are on a city's parks — a dog park
% survey and park lighting.
% Part A (the dog-park survey): A1 chance or not, A2 population and sample, A3 census and why
% not. Part B (parks and people): B1 the AP-style multiple choice (which plan is an experiment)
% with a required justification, B2 THE CRUX restated — step counters measured every step, so
% the paper calls it an experiment. 10 points.
% \answerspace{H}{} reserves exactly H; each \gradedblank inside it is BYTE-IDENTICAL in the key,
% where key mode prints its display. Keep every \boxguard identical in the blank and the key.
\documentclass[12pt]{article}
\usepackage{apstats-article}
\usepackage{apstats-key}   % pulls in -boxes; do NOT also load -boxes
\usepackage{grading}
\gradingassignment{apstats-u03-l01-hw}
\gradingcourse{advanced_statistics}
\newcommand{\answerspace}[2]{\par\nopagebreak\noindent\begin{minipage}[t][#1][t]{\linewidth}%
  \color{keyred}\bfseries #2\end{minipage}\par}

\begin{document}

\pageheader{Unit 3, Lesson 3.1}{Homework --- Answer Key}
{\small\itshape Write each final answer in its box. Only what is inside a box is graded.}\par
% NAMESTRIP: no \namedateperiod here — mirrors the blank.

\vspace{0.04in}

% Authored BYTE-IDENTICALLY in homework/ and homework_key/.
\boxguard[8]
\begin{remindbox}
\small
\textbf{This is your graded homework.} It is scored 1 to 5 and due the first class after two study
halls. Start it now; ask about anything that stalls you before you leave.
\end{remindbox}

\vspace{0.04in}

\boxguard[14]
\begin{notesbox}{Part A --- The Dog Park Survey}
\interlinepenalty=10000 % keep each prompt whole when the box breaks
\small
The city of Riverside has 40{,}000 households and wants to know what percent of them support
building a dog park. \textbf{Plan 1:} post a poll on the city's social media page and count the
votes. \textbf{Plan 2:} a computer picks 400 households at random from the city's list of all
addresses, and each one is asked.
\begin{enumerate}[label=\textbf{A\arabic*.}, leftmargin=2.2em, itemsep=2pt]
  \item Which plan uses chance to decide who answers? In the other plan, who decides?
        \answerspace{2.1in}{\normalcolor\mdseries\hspace*{10.5pt}\gradedblank[id=hw1, type=rubric, points=2, display={Plan 2: a computer picks the 400 households at random. In Plan 1 the residents decide, by choosing to see the post and vote.}]{\linewidth-10.5pt}{0.95cm}}
  \item For Plan 2, name the population and the sample.
        \answerspace{2.1in}{\normalcolor\mdseries\hspace*{10.5pt}\gradedblank[id=hw2, type=rubric, points=2, display={Population: all 40,000 Riverside households. Sample: the 400 households the computer picked.}]{\linewidth-10.5pt}{0.95cm}}
  \item A council member asks, ``Why not just ask every one of the 40{,}000 households?'' What is
        that plan called, and give one reason the city might not use it.
        \answerspace{2.1in}{\normalcolor\mdseries\hspace*{10.5pt}\gradedblank[id=hw3, type=rubric, points=2, display={A census. Reaching all 40,000 households would take far more time and money, and a random sample of 400 can estimate the percent well.}]{\linewidth-10.5pt}{0.95cm}}
\end{enumerate}
\end{notesbox}

\vspace{0.04in}

\boxguard[30]
\begin{notesbox}{Part B --- Parks and People}
\interlinepenalty=10000 % keep each prompt whole when the box breaks
\small
\begin{enumerate}[label=\textbf{B\arabic*.}, leftmargin=2.2em, itemsep=2pt]
  \item \textbf{AP-style.} The parks department wants to know whether adding lights to a park
        \emph{reduces} evening vandalism. Which plan is an experiment? Write its letter in the
        Letter box, then justify it in a sentence.
        \begin{enumerate}[label=\textbf{(\Alph*)}, leftmargin=2.2em, itemsep=0pt]
          \item Compare vandalism in parks that already have lights with parks that do not.
          \item Survey residents about whether lights would make parks safer.
          \item Ask park rangers to keep a careful daily log of all vandalism for a year.
          \item \ans{Randomly choose 10 of the city's 20 parks to get new lights, then compare
                vandalism in the two groups.}
          \item Review police reports from cities that chose to light their parks.
                \hfill\gradedright{\textit{Letter:}~\gradedblank[id=hw4, type=choice, points=1, answer=D]{0.45in}{0.24in}}\rule[-0.2in]{0pt}{0.48in}
        \end{enumerate}
        \answerspace{2.1in}{\normalcolor\mdseries\hspace*{10.5pt}\gradedblank[id=hw4j, type=rubric, points=1, display={(D): the department imposes the treatment (lights) on parks chosen at random. In the others, it only observes what already happens.}]{\linewidth-10.5pt}{0.7cm}}
  \item Researchers gave 600 Riverside adults step counters and recorded every step they took for
        a month. Dog owners walked more minutes a day than non-owners. The local paper wrote,
        ``Because every step was measured so carefully, this was an experiment --- getting a dog
        will make you walk more.'' Is the paper right? Explain in context.
        \answerspace{2.1in}{\normalcolor\mdseries\hspace*{10.5pt}\gradedblank[id=hw5, type=rubric, points=2, display={No. No one was assigned to get a dog, so it is an observational study, however carefully steps were measured. It shows an association, not that dogs cause walking; people who already like to walk may be the ones who get dogs.}]{\linewidth-10.5pt}{0.95cm}}
\end{enumerate}
\end{notesbox}

\vspace{0.04in}

% The next-lesson preview closes the packet, so it lives here rather than in AP Practice.
% A one-line, unbreakable box: guard it by its own height so it moves whole.
\boxguard[3]
\begin{spiralbox}
\small
\textbf{Next: Random Sampling Methods} --- how to let chance choose the sample: simple random
samples, stratified samples, and cluster samples.
\end{spiralbox}

\end{document}
